\chapter{AI Pipeline Design and Implementation}
\label{chap:contributions}

Chapter~\ref{chap:architecture} treated the recommendation engine and the hybrid catalogue mode as services with a fixed contract. This chapter looks inside both, then adds two shorter sections, the visual product search that fills the reference library and the print pipeline that prepares finished assets for paper.

The same idea runs through all three. The gap between what an off-the-shelf generative model produces and what a small retailer is willing to publish is rarely closed by a bigger model; it is closed by the protocol that sits around the model. The recommendation engine grounds the language model in retrieved context through a multi-stage prompt chain, the hybrid catalogue pipeline uses a deterministic compositor to bound what the image model is allowed to render, and the print pipeline tiles a small super-resolution model under fixed paper sizes.

\section{The Daily Recommendation Engine}
\label{sec:recommendation-engine}

\subsection{Problem statement and pipeline overview}
\label{subsec:rec-problem}

Small retailers know their products and their customers, but rarely have time to plan a daily marketing schedule. Out-of-the-box language models, asked to suggest a promotion for today, default to season-blind boilerplate: ``discover our premium selection'', ``don't miss out on our exclusive offer''. The recommendation engine's job is to produce, every day, a small ranked set of fresh marketing themes grounded in (a) the day's external context, (b) a curated promotional playbook of patterns that have worked in retail, and (c) the shop's own profile, without repeating what the shop has already seen in the recent past.

The engine's output is not a final asset, but more than just copy. Each suggestion is a structured record carrying a title, body copy, a suggested social-media caption, an asset \emph{category} (\emph{announcement} or \emph{promotion}) and an \texttt{imagePrompt} written specifically for the image-generation pipeline of Section~\ref{sec:hybrid-catalogue}. Tapping ``generate'' hands the retailer to the asset-creation screen with the prompt and category already filled in, so the path from a daily idea to a published asset is a single further confirmation. The engine is therefore the entry point of the platform's full creation flow, not a side feature. Figure~\ref{fig:rec-handoff} sketches that flow.

\begin{figure}[H]
\centering
% --- self-contained palette (muted, print-friendly) ---
\definecolor{cInk}{HTML}{2E3440}    % text + arrows
\definecolor{cIndigo}{HTML}{3F4D8C} % LLM / orchestration stages
\definecolor{cSlate}{HTML}{4A5568}  % deterministic platform code
\definecolor{cTeal}{HTML}{2E5F8C}   % stores
\definecolor{cAmber}{HTML}{B07A1F}  % highlight accent
\definecolor{cGray}{HTML}{7C8390}   % external providers
\resizebox{\textwidth}{!}{%
\begin{tikzpicture}[
    font=\sffamily\scriptsize,
    card/.style={draw=#1, line width=0.9pt, rounded corners=3pt, fill=#1!7,
                 align=center, text=cInk,
                 drop shadow={shadow xshift=0.5mm, shadow yshift=-0.5mm,
                              fill=black, opacity=0.12}},
    arr/.style={-{Latex[length=2mm,width=1.5mm]}, line width=0.7pt, draw=cInk!85},
    arrlbl/.style={font=\sffamily\tiny\itshape, text=cInk!80}
]

\node[card=cIndigo, minimum width=34mm, minimum height=20mm] (idea) at (0, 0) {%
    \textbf{Daily idea}\\[3pt]
    title, body, suggestedPost,\\
    category, imagePrompt
};
\node[font=\sffamily\tiny\itshape, text=cInk!70] at (0, -1.55) {(recommendation engine)};

\node[card=cSlate, minimum width=34mm, minimum height=20mm] (screen) at (5.0, 0)
    {Asset-creation\\screen\\[2pt]prompt and category\\pre-filled};

\node[card=cIndigo, minimum width=34mm, minimum height=20mm] (pipe) at (10.0, 0)
    {Generation\\pipeline\\[2pt](catalogue or\\announcement)};

\node[card=cTeal, minimum width=34mm, minimum height=20mm] (out) at (15.0, 0)
    {Final asset\\in gallery};

\draw[arr] (idea) -- node[above, arrlbl]{tap ``generate''} (screen);
\draw[arr] (screen) -- node[above, arrlbl]{confirm} (pipe);
\draw[arr] (pipe) -- (out);

\end{tikzpicture}%
}
\caption{Handoff from a daily idea to a published asset.}
\label{fig:rec-handoff}
\end{figure}

\label{subsec:rec-pipeline}
The engine has three parts that meet at inference time. An offline indexing path turns a curated corpus into the embeddings used at retrieval time. An offline fine-tuning path adapts a base writer model to the platform's voice. An online inference path stitches context signals, the shop profile and retrieved playbook entries into a prompt for the fine-tuned writer. Figure~\ref{fig:rec-pipeline} shows the arrangement at a glance.

\begin{figure}[H]
\centering
% --- self-contained palette (muted, print-friendly) ---
\definecolor{cInk}{HTML}{2E3440}    % text + arrows
\definecolor{cIndigo}{HTML}{3F4D8C} % LLM / orchestration stages
\definecolor{cSlate}{HTML}{4A5568}  % deterministic platform code
\definecolor{cTeal}{HTML}{2E5F8C}   % stores
\definecolor{cAmber}{HTML}{B07A1F}  % highlight accent
\definecolor{cGray}{HTML}{7C8390}   % external providers
\resizebox{\textwidth}{!}{%
\begin{tikzpicture}[
    font=\sffamily\scriptsize,
    card/.style={draw=#1, line width=0.9pt, rounded corners=3pt, fill=#1!7,
                 align=center, text=cInk,
                 drop shadow={shadow xshift=0.5mm, shadow yshift=-0.5mm,
                              fill=black, opacity=0.12}},
    group/.style={draw=#1, line width=0.9pt, rounded corners=3pt, fill=#1!6,
                  inner sep=3mm,
                  drop shadow={shadow xshift=0.5mm, shadow yshift=-0.5mm,
                               fill=black, opacity=0.12}},
    tab/.style={fill=#1, rounded corners=2.5pt, inner xsep=6pt, inner ysep=0pt,
                minimum height=5.2mm, anchor=south west,
                font=\sffamily\scriptsize\bfseries, text=white},
    chip/.style={draw=#1, line width=0.6pt, rounded corners=2.5pt, fill=white,
                 align=center, text=cInk},
    store/.style={draw=cTeal, line width=0.6pt, cylinder, shape border rotate=90,
                  aspect=0.22, minimum width=24mm, minimum height=15mm, align=center,
                  cylinder uses custom fill, cylinder body fill=white,
                  cylinder end fill=cTeal!15, text=cInk},
    person/.style={draw=cIndigo, line width=0.6pt, circle, fill=white,
                   minimum size=9mm, align=center, text=cInk},
    arr/.style={-{Latex[length=2mm,width=1.5mm]}, line width=0.7pt, draw=cInk!85},
    fb/.style={-{Latex[length=2mm,width=1.5mm]}, line width=0.7pt, draw=cGray,
               dash pattern=on 2.4pt off 1.7pt}
]

% --- Top: corpus
\node[card=cTeal, minimum width=34mm, minimum height=11mm] (data) at (0, 7.0)
    {Curated retail\\marketing corpus};

% --- Left branch: indexing
\node[chip=cSlate, minimum width=27mm, minimum height=8.5mm] (split) at (-3.8, 4.35)
    {split, chunk, embed};
\node[store] (vec) at (-3.8, 2.5) {pgvector\\(playbook +\\past ideas)};

% --- Right branch: fine-tuning
\node[chip=cIndigo, minimum width=30mm, minimum height=8.5mm] (td) at (3.8, 4.75)
    {training pairs};
\node[chip=cIndigo, minimum width=30mm, minimum height=11mm] (ft) at (3.8, 3.55)
    {LoRA fine-tune\\(Gemma 3 4B base)};
\node[chip=cIndigo, line width=0.9pt, fill=cIndigo!14, minimum width=30mm,
      minimum height=11mm] (ftllm) at (3.8, 2.1) {fine-tuned\\writer LLM};

\begin{pgfonlayer}{background}
\node[group=cSlate, fit=(split)(vec)] (idxbox) {};
\node[group=cIndigo, fit=(td)(ft)(ftllm)] (ftbox) {};
\end{pgfonlayer}

% --- Bottom: inference flow
\node[person] (retailer) at (-9.2, -1.5) {retailer};
\node[chip=cIndigo, minimum width=16mm, minimum height=8.5mm] (app) at (-7.3, -1.5)
    {app};

\node[chip=cAmber, minimum width=24mm, minimum height=11mm] (sig) at (-4.0, -1.5)
    {context signals};
\node[chip=cAmber, minimum width=20mm, minimum height=11mm] (req) at (-1.55, -1.5)
    {shop profile};
\node[chip=cAmber, minimum width=31mm, minimum height=11mm] (ctx) at (1.55, -1.5)
    {retrieved playbook\\+ past ideas};

\begin{pgfonlayer}{background}
\node[group=cAmber, inner sep=2.8mm, fit=(sig)(req)(ctx)] (ragbox) {};
\end{pgfonlayer}

\node[card=cIndigo, minimum width=20mm, minimum height=10mm] (out) at (6.9, -1.5)
    {ideas};

% --- Arrows: corpus splits into both branches
\draw[arr] (data.south) to[out=-115, in=20] (split.east);
\draw[arr] (data.south) to[out=-65, in=160] (td.west);

% Indexing internal
\draw[arr] (split) -- (vec);

% Fine-tuning internal
\draw[arr] (td) -- (ft);
\draw[arr] (ft) -- (ftllm);

% Inference flow
\draw[arr] (retailer) -- (app);
\draw[arr] (app) -- (ragbox.west);
\draw[arr] (vec.south) to[out=-90, in=90] (ctx.north);
\draw[arr] (ragbox.east) to[out=0, in=-90] (ftllm.south);
\draw[arr] (ftllm.east) to[out=0, in=90] (out.north);

% Feedback loop
\draw[fb] (out.south) to[out=-90, in=-90, looseness=0.4]
  node[below=1pt, font=\sffamily\tiny\itshape, text=cGray, sloped]{retailer feedback}
  (retailer.south);

% --- Group tabs (drawn last so they sit above the arrows)
\node[tab=cSlate]  at ([xshift=4mm]idxbox.north west) {Indexing};
\node[tab=cIndigo] at ([xshift=4mm]ftbox.north west)  {Fine-tuning};
\node[tab=cAmber]  at ([xshift=4mm]ragbox.north west) {RAG prompt};

\end{tikzpicture}%
}
\caption{Architecture of the recommendation engine.}
\label{fig:rec-pipeline}
\end{figure}

The inference path itself is organised as a six-stage chain. The stages are, in order: (1) context aggregation (Section~\ref{subsec:rec-context}), (2) playbook retrieval (Section~\ref{subsec:rec-playbook}), (3) anti-repetition retrieval (Section~\ref{subsec:rec-anti-repetition}), (4) the strategist, (5) the writers and (6) the translator, with stages 4 to 6 covered together in Section~\ref{subsec:rec-llm-chain}. Stages 1 to 3 produce the artefacts that ground the prompt, and stages 4 to 6 are the language-model calls that turn those artefacts into ideas. Each stage produces a small, well-typed artefact that is consumed by the next, and a failure in any single stage degrades the run gracefully rather than aborting it. The choice of six stages, rather than a single end-to-end prompt, is deliberate. Each stage has a narrowly scoped responsibility, exchanges a typed artefact rather than free text with its neighbours, and can be observed independently. When an idea is judged poor, the per-stage trace reveals whether the failure was in the signal, the retrieval, the planning or the drafting, which would not be possible if the whole pipeline were a single opaque call. The remaining subsections walk through the six stages in that order.

\subsection{External context signals and anti-repetition}
\label{subsec:rec-context}
\label{subsec:rec-anti-repetition}

The engine currently consumes three classes of external signal, with room to add more later. Each runs in parallel and is wrapped in its own error boundary, so a single failure does not abort the run. Failed sources are recorded and the rest continue.

\paragraph{Weather.} A short-range forecast detects heatwaves, cold snaps, heavy rain and storms, and turns each event into a signal tagged high or medium with a one-line summary the prompt can read directly.

\paragraph{Holidays.} Upcoming public holidays in the shop's country are emitted over a fourteen-day lookahead, with severity scaling by proximity: high at most three days away, medium at most seven, low further out.

\paragraph{News.} Recent news items in the shop's country are emitted at low severity as ambient context, useful for seeding an angle but never for overriding the day's plan.

\begin{table}[!ht]
\centering
\caption{Context-signal severity rules.}
\label{tab:rec-signals}
\small
\renewcommand{\arraystretch}{1.3}
\begin{tabularx}{\textwidth}{|>{\raggedright\arraybackslash}p{2.4cm}|Y|}
\hline
\textbf{Signal} & \textbf{Severity rule} \\
\hline
Weather  & High for heatwave, cold snap or storm, medium for heavy-rain probability $\geq 70\,\%$ without a storm. \\
\hline
Holiday  & High at $\leq 3$ days, medium at $\leq 7$, low further out. \\
\hline
News     & Always low, ambient context only. \\
\hline
\end{tabularx}
\end{table}

A complementary signal, used at retrieval time rather than at prompt time, is the engine's own recent output. A second pgvector index records the embeddings of every idea the engine has ever produced for a tenant. After each run, the title and body of each new idea are concatenated, embedded by the same encoder, and stored in the \texttt{IdeaEmbedding} table introduced in Section~\ref{subsec:relational-schema}. At the start of the next run, the engine retrieves the five nearest neighbours of the current run's retrieval query within a thirty-day window and injects them, prefixed by a strict ``do not repeat these themes'' clause, into the writer prompt.

The arrangement converts memory into retrieval: the engine has no notion of ``what was suggested last week'' beyond what the embedding space surfaces, but the embedding space surfaces semantically close ideas regardless of the exact words used to express them. A retailer who saw a ``rainy-day comfort food'' theme on Tuesday is not shown a ``cosy-weather snack box'' theme on Friday, even though no string match would catch the duplication.

\subsection{The promotional playbook and retrieval-augmented generation}
\label{subsec:rec-playbook}

The engine's second grounding source is the platform's \emph{promotional playbook}: a curated table of marketing patterns that have been observed to work in independent retail. Each entry carries a \texttt{Theme} (e.g., ``rainy-day comfort food''), a \texttt{TriggerType} (\emph{weather}, \emph{holiday}, \emph{news} or \emph{seasonal}), a free-text \texttt{Trigger} description, a list of \texttt{Tactics} and an \texttt{ExampleCopy} excerpt, alongside a text embedding of those fields used at retrieval time.

The playbook is partitioned by business domain. The platform ships with a dedicated playbook for each of the four named domains it supports. Market covers groceries, household goods, snacks and drinks. Food covers caf\'es, bakeries and prepared meals. Retail covers general-goods shops. Fashion covers clothing and accessories. A fifth Other playbook of domain-agnostic patterns takes over for shops whose domain falls outside the four. Adding a new named domain is purely an authoring exercise. A fresh set of entries is ingested under the same schema, no code change is required, and the engine picks the right playbook by matching the shop's recorded domain.

The entries themselves were not generated by a model. They were authored by hand from three sources: marketing literature aimed at independent retailers, a manual review of the social-media output of small shops in the target region, and observations from the development team's pilot outreach. A pattern was admitted to the playbook only when it recurred across all three sources, so every entry corresponds to a tactic with a known precedent rather than to a novel suggestion the engine has to invent on its own. The result is a small corpus, in the low hundreds of entries, that is dense, opinionated and easy to inspect, which is exactly the regime in which retrieval-augmented prompting outperforms an end-to-end model call.

At run-time the engine builds a short retrieval query from the top-ranked context signals and the shop's profile (for example: \emph{``Market shop ``Pet Store'' \(\mid\) weather: heatwave starting May 12 \(\mid\) holiday: Mother's Day in 2 days''}), embeds it through the same encoder used at ingestion, and retrieves the three nearest entries by cosine similarity within the playbook that matches the shop's domain, falling back to the general-purpose playbook when no domain-specific match is found. The retrieval is implemented over an HNSW index \cite{malkov2018hnsw} on the pgvector column, and the top-three entries are injected verbatim into the strategist prompt of Section~\ref{subsec:rec-llm-chain}. This is the retrieval-augmented generation pattern of Lewis et al.~\cite{lewis2020rag} adapted to a small, hand-curated corpus rather than to a large open knowledge base, on the principle that the marketing patterns small retailers benefit from are few, well known to specialists and best authored once than re-derived per request.

\subsection{The multi-stage language-model chain}
\label{subsec:rec-llm-chain}

The fourth, fifth and sixth stages are language-model calls organised as a chain rather than a single end-to-end prompt. The split lets each stage run at the temperature, JSON contract and model size that suit its task, and lets the writer stage fan out in parallel.

\paragraph{Stage 4: the Strategist.}
A single language-model call, at temperature $0.4$, is given the shop profile, the top context signals, the three retrieved playbook entries and a brief instruction to produce $N$ distinct angles. Its output is a JSON document of the form
\[
\{\text{angles}: [\,\{\text{theme},\ \text{tone},\ \text{triggerSignal},\ \text{rationale}\},\ \dots\,]\}.
\]
The low temperature is appropriate here because planning benefits from consistency: the strategist's job is to spread the angles across the available signals, not to be inventive on each one.

\paragraph{Stage 5: the Writers.}
For each angle the strategist produced, a separate language-model call is issued in parallel, at temperature $0.8$. The writer is given the angle, the shop identity, and the anti-repetition block from Section~\ref{subsec:rec-anti-repetition}. Its output is a JSON document
\[
\{\text{id},\ \text{tone},\ \text{title},\ \text{meta},\ \text{body},\ \text{suggestedPost},\ \text{type},\ \text{imagePrompt}\}.
\]
The high temperature is appropriate here because drafting benefits from variation: identical angles drafted at low temperature produce ideas that are too uniform across days. When the retailer accepts a recommendation, the \texttt{type} (\emph{announcement} or \emph{promotion}) and \texttt{imagePrompt} fields redirect them to the generation studio on the matching asset flow (see the \emph{Asset generation} requirement in Section~\ref{subsec:functional-requirements}, p.~\pageref{par:asset-generation}) with the prompt already filled in.

\paragraph{Stage 6: the Translator and persistence.}
For platforms set to a non-English interface language, each English idea is translated by a parallel language-model call into the target language. The translator prompt explicitly forbids a list of marketing clich\'es (in the Romanian case: ``descoper\u{a}'', ``Premium'', generic discount phrasing) so that the localised output does not regress to the boilerplate that the playbook and writer stages were specifically chosen to avoid. The translated idea is stored alongside the English one on the same record, and the gallery picks the right language at read time. Persistence then writes the day's run into the \texttt{DailyRecommendation} table together with a context snapshot, after which Stage~6 issues the embedding writes that fuel the next day's anti-repetition retrieval.

Table~\ref{tab:rec-stages} summarises the per-stage parameters.

\begin{table}[!ht]
\centering
\caption{Per-stage parameters of the recommendation chain.}
\label{tab:rec-stages}
\small
\renewcommand{\arraystretch}{1.3}
\begin{tabularx}{\textwidth}{|>{\raggedright\arraybackslash}p{2.6cm}|>{\raggedright\arraybackslash}p{2.0cm}|>{\raggedright\arraybackslash}p{1.5cm}|Y|}
\hline
\textbf{Stage} & \textbf{Calls per run} & \textbf{Temp.} & \textbf{Output shape} \\
\hline
Strategist  & 1                       & 0.4 & \texttt{\{angles: [\{theme, tone, triggerSignal, rationale\}\,\dots\,]\}} \\
\hline
Writer      & one per angle, parallel & 0.8 & \texttt{\{title, meta, body, suggestedPost, imagePrompt, type\}} \\
\hline
Translator  & one per idea, parallel  & 0.8 & translated counterpart of the writer output \\
\hline
\end{tabularx}
\end{table}

\subsection{Model selection and writer fine-tuning}
\label{subsec:rec-model-selection}

Several language models were trialled against the same fixed batch of $(\text{shop}, \text{day}, \text{signal})$ inputs and compared on three axes: tonal fit, specificity to the \texttt{triggerSignal}, and JSON-schema adherence. The candidates were Gemma 3 4B, DeepSeek, a LoRA-fine-tuned Gemma 3 4B writer, and the managed APIs Claude Haiku 4.5 and GPT-4o mini. No single model was best on every axis. All of the candidates kept their output on topic, but they varied a lot on brand voice and on how closely they followed the schema. Claude Haiku 4.5 read the most natural and stayed closest to the schema, while GPT-4o mini was looser on both. The plain Gemma 3 4B model was the weakest on tone of voice even though it kept up on the other two axes, so brand voice, rather than basic competence, is its real gap. The fine-tuned Gemma 3 4B writer improved on the plain model's voice at a much lower hosting cost, but it did not yet match the managed models on that axis. Table~\ref{tab:rec-model-comparison} summarises the comparison, and its ratings line up with the per-axis DeepEval scores reported in Chapter~\ref{chap:results-testing}. Closing the voice gap is the immediate motivation for the writer-stage fine-tuning described next.

{\small\renewcommand{\arraystretch}{1.3}
\begin{tabularx}{\textwidth}{|>{\raggedright\arraybackslash}p{2.8cm}|Y|Y|Y|Y|}
\caption{Qualitative comparison of the candidate models on the three evaluation axes.}\label{tab:rec-model-comparison}\\
\hline
\textbf{Model} & \textbf{Tonal fit} & \textbf{Specificity} & \textbf{JSON adherence} & \textbf{Hosting cost} \\
\hline
\endfirsthead

\multicolumn{5}{l}{\textit{Table~\ref{tab:rec-model-comparison} continued from previous page}}\\
\hline
\textbf{Model} & \textbf{Tonal fit} & \textbf{Specificity} & \textbf{JSON adherence} & \textbf{Hosting cost} \\
\hline
\endhead

\hline
\endfoot

\hline
\endlastfoot

Claude Haiku 4.5        & Good      & Strong    & Good      & Medium \\
\hline
GPT-4o mini             & Fair      & Strong    & Fair      & Medium \\
\hline
Fine-tuned Gemma 3 4B (LoRA) & Fair      & Strong    & Good      & Very low \\
\hline
DeepSeek                & Good      & Excellent & Fair      & Low  \\
\hline
Gemma 3 4B              & Weak      & Good      & Good      & Very low \\
\hline
\end{tabularx}}

\label{subsec:rec-finetune}
The strategist and translator stages are well served by an off-the-shelf instruction-tuned base model, because their tasks (producing a small structured plan and translating short text) are close enough to the model's pre-training distribution that prompt engineering is sufficient. The writer stage is the place where this is no longer true. A base Gemma 3 4B instruction model, even when given the strategist angle, the playbook entries and the anti-repetition block, gravitates back to the same handful of marketing tropes (``Discover our premium selection'', ``Don't miss out''), regardless of how strongly the prompt forbids them.

The platform therefore serves the writer role with a fine-tuned variant of the Gemma 3 4B instruction model, adapted via low-rank adaptation (LoRA) \cite{hu2022lora} on a curated corpus of marketing copy from independent retailers. LoRA was preferred over full fine-tuning because the corpus is small, the base model is already well aligned to instruction following, and the low-rank adapters can be hot-swapped at deployment time without redistributing the base weights. The corpus pairs context-signal triples with handcrafted copy that exhibits the qualities the boilerplate base model lacks: a specific, time-bounded reaction to the signal, a voice consistent with a single shop owner, and a concrete call to action that names a real product or aisle rather than a generic category. The strategist and translator continue to use the un-adapted base model.

The choice to fine-tune only the writer is itself part of the protocol-over-model-size argument that runs through this chapter: the writer is the one stage where a base model's failures are systematic and visible to the retailer, and the cheapest fix is to adapt the model to the platform's voice rather than ask the prompt to do all the work.

\section{Image Generation Pipelines}
\label{sec:hybrid-catalogue}

\subsection{Asset paths and prompting strategy}
\label{subsec:gen-overview}

The platform produces three families of marketing assets: announcements, wallpapers and catalogues. Each family has its own service, but the three services share the same back end. Each service assembles a prompt and an optional list of inline reference images, then hands both to the \texttt{IImageProvider} interface, which forwards them to whichever provider the retailer picked at request time (the Gemini provider in production, a mock provider in tests). What differs between services is only the prompt template and which inputs become inline images. Figure~\ref{fig:image-provider-interface} sketches the arrangement.

A single-shot prompt covers all three. Generation runs from a phone, has to finish in a few seconds, and is paid for in one inference call. Because the platform already knows the asset family before building the prompt, it can pre-select the right template and fill in the retailer's brand context, format and reference photos directly.

\begin{figure}[H]
\centering
% --- self-contained palette (muted, print-friendly) ---
\definecolor{cInk}{HTML}{2E3440}    % text + arrows
\definecolor{cIndigo}{HTML}{3F4D8C} % LLM / orchestration stages
\definecolor{cSlate}{HTML}{4A5568}  % deterministic platform code
\definecolor{cTeal}{HTML}{2E5F8C}   % stores
\definecolor{cAmber}{HTML}{B07A1F}  % highlight accent
\definecolor{cGray}{HTML}{7C8390}   % external providers
\resizebox{0.85\textwidth}{!}{%
\begin{tikzpicture}[
    font=\sffamily\scriptsize,
    svc/.style={draw=cSlate, line width=0.9pt, rounded corners=3pt, fill=cSlate!7,
                minimum width=32mm, minimum height=10mm, align=center, text=cInk,
                drop shadow={shadow xshift=0.5mm, shadow yshift=-0.5mm,
                             fill=black, opacity=0.12}},
    iface/.style={draw=cIndigo, line width=1.0pt, rounded corners=3pt,
                  fill=cIndigo!14, minimum width=50mm, minimum height=18mm,
                  align=center, text=cInk,
                  drop shadow={shadow xshift=0.5mm, shadow yshift=-0.5mm,
                               fill=black, opacity=0.12}},
    prov/.style={draw=cGray, line width=0.9pt, rounded corners=3pt, fill=cGray!6,
                 dash pattern=on 2.4pt off 1.7pt, align=center, text=cInk},
    group/.style={draw=#1, line width=0.9pt, rounded corners=3pt, fill=#1!6,
                  inner sep=2.8mm,
                  drop shadow={shadow xshift=0.5mm, shadow yshift=-0.5mm,
                               fill=black, opacity=0.12}},
    tab/.style={fill=#1, rounded corners=2.5pt, inner xsep=6pt, inner ysep=0pt,
                minimum height=5.2mm, anchor=south west,
                font=\sffamily\scriptsize\bfseries, text=white},
    bus/.style={line width=0.9pt, draw=cInk!70, rounded corners=2pt},
    arr/.style={-{Latex[length=2mm,width=1.5mm]}, line width=0.7pt, draw=cInk!85},
    arrlbl/.style={font=\sffamily\tiny\itshape, text=cInk!80}
]

% --- asset services (deterministic platform code) ---
\node[svc] (ann) at (0, 1.7)  {Announcement\\service};
\node[svc] (wal) at (0, 0)    {Wallpaper\\service};
\node[svc] (cat) at (0, -1.7) {Catalogue\\service};

\begin{pgfonlayer}{background}
\node[group=cSlate, fit=(ann)(wal)(cat)] (svcbox) {};
\end{pgfonlayer}

% --- the shared abstraction (hero) ---
\node[iface] (iface) at (6.9, 0)
    {\textbf{\texttt{IImageProvider}}\\[3pt]
     {\color{cIndigo}\texttt{GenerateAsync(}}\\
     {\color{cIndigo}\texttt{prompt, inlineImages)}}};
\node[arrlbl, text=cInk!70] at (6.9, -1.42) {single seam, provider-agnostic};

% --- external providers ---
\node[prov, minimum width=40mm, minimum height=18mm] (prov) at (13.8, 0)
    {Image generation providers\\(Gemini, OpenAI, \dots)};

% --- three callers merge onto one bus, one shared call ---
\draw[bus] (ann.east) -- (2.7, 1.7) -- (2.7, -1.7) -- (cat.east);
\draw[bus] (wal.east) -- (2.7, 0);
\fill[cInk!70] (2.7, 0) circle (1.6pt);
\draw[arr] (2.7, 0) -- node[above=2pt, arrlbl]{shared call} (iface.west);

% --- provider chosen per request ---
\draw[arr] (iface.east) -- node[above=2pt, arrlbl]{selected per request} (prov.west);

% --- group tabs (drawn last, above the arrows) ---
\node[tab=cSlate] at ([xshift=4mm]svcbox.north west) {Asset services};
\node[tab=cGray]  at ([xshift=4mm]prov.north west)   {External};

\end{tikzpicture}%
}
\caption{All three asset services share the same \texttt{IImageProvider} call, with the provider chosen per request.}
\label{fig:image-provider-interface}
\end{figure}

\textbf{Announcement.} Each of the three sub-types listed in the functional requirements is rendered by its own prompt template behind the same entry point. The general announcement template hands the model the raw content and lets it pick between an event headline and a tip-card layout. The job-post template assembles the prompt from structured fields (title, schedule, salary, requirements, with-person or text-only style) and applies a no-invented-CTA rule so application instructions appear only when the retailer typed them in. The promotion template treats the discount as the hero and folds in product reference photos when supplied. A single-shot prompt fits all three because the variation lives in the inputs, not in the reasoning.

\textbf{Wallpaper.} Wallpapers are background templates onto which real product photos are pasted later. The prompt is built dynamically from the brand context the retailer has filled in. A header band is reserved only when there is brand text or a logo, a footer band only when there is contact or social information, and otherwise the whole canvas is given over to a clean hero zone. Two rules close every wallpaper prompt: the hero area must stay plain, with no patterns or props, and the only text the model may render is the exact strings the prompt lists.

\textbf{Catalogue.} The catalogue path turns the retailer's product photos, brand context, layout style, colour theme and format into a single catalogue image. The prompt carries an art-direction block (composition, typography, lighting, colour) so that one call to the provider produces a usable result without a separate critique step. Catalogue generation has two model-based modes, alongside a third path that invokes no generative model at all and is described in Section~\ref{subsec:inhouse-composite}. The default mode lets the model render everything, the products included. The preserve mode keeps the retailer's product pixels untouched.

The preserve mode is the most complex path on the platform. It is the only one that combines a generative step with a deterministic post-processor, and the only one that imposes a strict contract on the model output, namely a coloured silhouette around each product and a quiet zone of clean background around each silhouette, that the post-processor then relies on. It exists because some retailers cannot tolerate any change to their products, even small rewrites of label text or substitutions with visually similar items, which off-the-shelf multimodal models still produce regularly. The rest of this section explores how the preserve mode is built, why a fully generative pipeline does not meet that bar, and how the deterministic compositor closes the gap.

\subsection{Why fully generative is insufficient and the outline-as-protocol idea}
\label{subsec:hyb-motivation}

A retail catalogue must show the actual products on offer. A photograph of a packet of crisps that subtly differs from the real packet, in label colour, in label text or in shape, is not a marketing asset; it is a misrepresentation of the goods, and a small retailer cannot publish it without risking customer trust and, in some jurisdictions, falling foul of advertising law. The bar for product fidelity is therefore much higher than for the surrounding scene.

Even the best multimodal image models, when conditioned on the retailer's own product photographs as inline references, fail this bar with regularity. Two failure modes recur. The first is \emph{label hallucination}: the model preserves the rough shape and dominant colour of the product but invents the text on its label, producing a packet whose typeface, brand mark and ingredient list are entirely fictional. The second is \emph{lookalike substitution}: the model replaces the reference product with a different but visually similar one, for example a near-identical chocolate bar from a competing brand. A small but persistent failure rate of either kind is unacceptable for marketing material that goes out under the retailer's name.

The hybrid catalogue mode introduced in Section~\ref{sec:requirements} of Chapter~\ref{chap:architecture} is the platform's response to this constraint. Its central idea is to decouple scene generation, which generative models do well, from product fidelity, which they do badly, and to enforce the latter with a small deterministic post-processor.

\label{subsec:hyb-protocol}
The platform's response to the constraint above is the outline-as-protocol idea. The platform asks the model to render the full catalogue scene with the retailer's products in place, and to draw a tight silhouette outline around each product in a unique marker colour reserved for that product. The model is told to keep every product faithful to its reference in size, shape and orientation, which current models do reliably. A deterministic compositor then detects each marker colour, erases the rendered product together with its outline, inpaints the residue, and pastes the retailer's real product photograph into the cleaned region at the same position and size. The product is therefore untouched, pixel for pixel, while the model handles the backdrop, props, colour grading, layout and accompanying text.

The mode exists for retailers who cannot tolerate any change to their packaging. Even visually similar renders often rewrite label text, alter typography or invent small graphical details that the retailer spots at a glance. The hybrid path leaves the packaging exactly as the retailer supplied it.

\begin{figure}[H]
\centering
% --- self-contained palette (muted, print-friendly) ---
\definecolor{cInk}{HTML}{2E3440}    % text + arrows
\definecolor{cIndigo}{HTML}{3F4D8C} % LLM / orchestration stages
\definecolor{cSlate}{HTML}{4A5568}  % deterministic platform code
\definecolor{cTeal}{HTML}{2E5F8C}   % stores
\definecolor{cAmber}{HTML}{B07A1F}  % highlight accent
\definecolor{cGray}{HTML}{7C8390}   % external providers
\resizebox{\textwidth}{!}{%
\begin{tikzpicture}[
    font=\sffamily\scriptsize,
    card/.style={draw=#1, line width=0.9pt, rounded corners=3pt, fill=#1!7,
                 minimum width=30mm, minimum height=12mm, align=center, text=cInk,
                 drop shadow={shadow xshift=0.5mm, shadow yshift=-0.5mm,
                              fill=black, opacity=0.12}},
    model/.style={draw=cGray, line width=0.9pt, rounded corners=3pt, fill=cGray!6,
                  dash pattern=on 2.4pt off 1.7pt, minimum width=30mm,
                  minimum height=12mm, align=center, text=cInk},
    ghost/.style={draw=cGray!55, line width=0.8pt, rounded corners=3pt, fill=cGray!4,
                  dash pattern=on 2pt off 1.6pt, minimum width=30mm,
                  minimum height=12mm, align=center, text=cGray!85,
                  font=\sffamily\scriptsize\itshape},
    group/.style={draw=#1, line width=0.9pt, rounded corners=3pt, fill=#1!6,
                  inner sep=2.8mm,
                  drop shadow={shadow xshift=0.5mm, shadow yshift=-0.5mm,
                               fill=black, opacity=0.12}},
    tab/.style={fill=#1, rounded corners=2.5pt, inner xsep=6pt, inner ysep=0pt,
                minimum height=5.2mm, anchor=south west,
                font=\sffamily\scriptsize\bfseries, text=white},
    arr/.style={-{Latex[length=2mm,width=1.5mm]}, line width=0.7pt, draw=cInk!85},
    darr/.style={-{Latex[length=2mm,width=1.5mm]}, line width=0.7pt, draw=cInk!70,
                 dash pattern=on 2.4pt off 1.7pt},
    arrlbl/.style={font=\sffamily\tiny\itshape, fill=white, inner sep=1pt,
                   text=cInk!80}
]

% --- (a) pure generative: no guardrail stage, products distort ---
\node[card=cTeal]  (refA) at (0,    2.7) {Reference\\photographs};
\node[model]       (mdlA) at (4.4,  2.7) {Multimodal\\model};
\node[ghost]       (ghA)  at (9.2,  2.7) {no compositor};
\node[card=cAmber] (outA) at (13.8, 2.7) {Final image\\\textit{products distorted}};
\draw[arr] (refA) -- (mdlA);
\draw[arr] (mdlA) -- (ghA);
\draw[arr] (ghA)  -- (outA);

% --- (b) hybrid: outline mode + deterministic compositor (the contribution) ---
\node[card=cTeal]  (refB) at (0,    0) {Reference\\photographs};
\node[model]       (mdlB) at (4.4,  0) {Multimodal model\\(outline mode)};
\node[card=cSlate, line width=1.0pt, fill=cSlate!16] (cmpB) at (9.2, 0)
    {Deterministic\\compositor};
\node[card=cTeal]  (outB) at (13.8, 0) {Final image\\\textit{products faithful}};
\draw[arr] (refB) -- (mdlB);
\draw[arr] (mdlB) -- node[arrlbl, above=1pt]{outlined scene} (cmpB);
\draw[arr] (cmpB) -- (outB);

% --- real reference photographs also feed the compositor ---
\draw[darr] (refB.south)
    to[out=-55, in=-125, looseness=0.9]
    node[arrlbl, below=2pt, name=fblbl]{real product photos}
    (cmpB.south);

% --- framed rows with tabs (drawn behind the nodes) ---
\begin{pgfonlayer}{background}
\node[group=cGray,  fit=(refA)(mdlA)(ghA)(outA)]            (rowA) {};
\node[group=cSlate, fit=(refB)(mdlB)(cmpB)(outB)(fblbl)]    (rowB) {};
\end{pgfonlayer}
\node[tab=cGray]  at ([xshift=4mm]rowA.north west) {(a) Pure generative};
\node[tab=cSlate] at ([xshift=4mm]rowB.north west) {(b) Hybrid: outline + composite};

\end{tikzpicture}%
}
\caption{Pure generative path versus hybrid path}
\label{fig:hybrid-vs-pure}
\end{figure}

\subsection{The three-stage pipeline}
\label{subsec:hyb-pipeline}

\begin{figure}[H]
\centering
% --- self-contained palette (muted, print-friendly) ---
\definecolor{cInk}{HTML}{2E3440}    % text + arrows
\definecolor{cIndigo}{HTML}{3F4D8C} % LLM / orchestration stages
\definecolor{cSlate}{HTML}{4A5568}  % deterministic platform code
\definecolor{cTeal}{HTML}{2E5F8C}   % stores
\definecolor{cAmber}{HTML}{B07A1F}  % highlight accent
\definecolor{cGray}{HTML}{7C8390}   % external providers
\resizebox{\textwidth}{!}{%
\begin{tikzpicture}[
    font=\sffamily\scriptsize,
    group/.style={draw=#1, line width=0.9pt, rounded corners=3pt, fill=#1!6,
                  inner sep=3mm,
                  drop shadow={shadow xshift=0.5mm, shadow yshift=-0.5mm,
                               fill=black, opacity=0.12}},
    tab/.style={fill=#1, rounded corners=2.5pt, inner xsep=6pt, inner ysep=0pt,
                minimum height=5.2mm, anchor=south west,
                font=\sffamily\scriptsize\bfseries, text=white},
    inp/.style={draw=#1, line width=0.6pt, rounded corners=2.5pt, fill=white,
                minimum width=28mm, minimum height=8mm, align=center, text=cInk,
                font=\sffamily\scriptsize\itshape},
    proc/.style={draw=#1, line width=0.6pt, rounded corners=2.5pt, fill=white,
                 minimum width=32mm, minimum height=11mm, align=center, text=cInk},
    aicall/.style={draw=cGray, line width=0.9pt, rounded corners=3pt, fill=cGray!6,
                   dash pattern=on 2.4pt off 1.7pt, minimum width=34mm,
                   minimum height=13mm, align=center, text=cInk},
    art/.style={draw=cAmber, line width=0.6pt, rounded corners=2.5pt, fill=white,
                minimum width=30mm, minimum height=9mm, align=center, text=cInk,
                font=\sffamily\scriptsize\itshape},
    final/.style={draw=cTeal, line width=0.9pt, rounded corners=3pt, fill=cTeal!7,
                  minimum width=32mm, minimum height=10mm, align=center, text=cInk,
                  drop shadow={shadow xshift=0.5mm, shadow yshift=-0.5mm,
                               fill=black, opacity=0.12}},
    arr/.style={-{Latex[length=2mm,width=1.5mm]}, line width=0.7pt, draw=cInk!85}
]

% ============ Stage 1: Marker palette ============
\node[inp=cSlate]  (in1a)       at (0, 5.5) {brand colours};
\node[inp=cSlate]  (in1b)       at (0, 4.4) {layout theme};
\node[inp=cSlate]  (in1c)       at (0, 3.3) {products};

\node[proc=cSlate] (palettesel) at (3.8, 4.4) {Marker palette\\selection};
\node[art]         (palette)    at (3.8, 2.2) {marker assignment\\$\textit{product} \rightarrow \textit{hex}$};

\draw[arr] (in1a.east) -- (palettesel.west);
\draw[arr] (in1b.east) -- (palettesel.west);
\draw[arr] (in1c.east) -- (palettesel.west);
\draw[arr] (palettesel.south) -- (palette.north);

\begin{pgfonlayer}{background}
\node[group=cSlate, fit=(in1a)(in1c)(palettesel)(palette)] (stage1) {};
\end{pgfonlayer}

% ============ Stage 2: Outline rendering ============
\node[inp=cIndigo, minimum width=24mm] (in2a) at (7.4, 5.5) {scene prompt};
\node[inp=cIndigo, minimum width=27mm] (in2b) at (10.2, 5.5) {reference photos};

\node[aicall] (mdl)    at (8.8, 3.2) {Multimodal model\\(outline mode)\\[1pt]\tiny\textsc{ai call}};
\node[art, minimum width=34mm] (canvas) at (8.8, 0.9)
    {placeholder canvas\\(scene + outlined products)};

\draw[arr] (in2a.south) -- (mdl.north);
\draw[arr] (in2b.south) -- (mdl.north);
\draw[arr] (palette.east) to[out=0, in=180] (mdl.west);
\draw[arr] (mdl.south) -- (canvas.north);

\begin{pgfonlayer}{background}
\node[group=cIndigo, fit=(in2a)(in2b)(mdl)(canvas)] (stage2) {};
\end{pgfonlayer}

% ============ Stage 3: Detect and composite ============
\node[proc=cSlate, minimum width=38mm] (detect) at (14.6, 5.3)
    {Marker detect\\(Chebyshev match,\\largest component)};
\node[proc=cSlate, minimum width=38mm] (mask)   at (14.6, 3.6)
    {Mask construction\\(flood-fill + dilate)};
\node[aicall, minimum width=38mm, minimum height=15mm] (lama) at (14.6, 1.7)
    {LaMa inpainting\\(erase outlines,\\repaint background)\\[1pt]\tiny\textsc{ai call (IOPaint)}};
\node[proc=cSlate, minimum width=38mm] (paste)  at (14.6, -0.3)
    {Scale \& alpha-composite\\(retailer photo paste)};

\draw[arr] (canvas.east) to[out=0, in=180, looseness=0.8] (detect.west);
\draw[arr] (detect.south) -- (mask.north);
\draw[arr] (mask.south)   -- (lama.north);
\draw[arr] (lama.south)   -- (paste.north);

\node[final] (out) at (14.6, -2.2) {Final image\\(products faithful)};
\draw[arr] (paste.south) -- (out.north);

\begin{pgfonlayer}{background}
\node[group=cSlate, fit=(detect)(mask)(lama)(paste)] (stage3) {};
\end{pgfonlayer}

% ============ Stage tabs (drawn last so they sit above the arrows) ============
\node[tab=cSlate]  at ([xshift=4mm]stage1.north west) {Stage 1: marker palette};
\node[tab=cIndigo] at ([xshift=4mm]stage2.north west) {Stage 2: outline rendering};
\node[tab=cSlate]  at ([xshift=4mm]stage3.north west) {Stage 3: detect and composite};

\end{tikzpicture}%
}
\caption{Three-stage pipeline of the hybrid catalogue mode: marker palette selection, outline rendering by the multimodal model, and deterministic detect-erase-composite.}
\label{fig:hybrid-pipeline}
\end{figure}

The pipeline has three stages: a marker-palette selection step that decides which colours to reserve for which products, an outline-rendering step in which the multimodal model produces a placeholder catalogue, and a deterministic detect-and-composite step that turns the placeholder into the final image. Figure~\ref{fig:hybrid-pipeline} shows the artefacts that flow between the stages.

\paragraph{Stage 1: marker palette selection.}
Before any model call is issued, the platform decides which colour will mark which product. The candidate pool is twelve highly saturated neon hues (magenta, cyan, electric violet, hot pink, chartreuse and so on), ordered so that the first $k$ entries are spread evenly around the hue circle for any $k$ from one to twelve. From this pool two exclusion rules remove unsuitable colours:

\begin{enumerate}
\item Brand-aware exclusion: candidates within Euclidean RGB distance $80$ of any brand-palette colour are dropped, so an outline cannot collide with a brand-coloured scene element at detection time.
\item Theme-aware exclusion: each catalogue layout theme carries a small set of colours that the multimodal model tends to emit as scene fill (for example, in a Pop-Art or Vibrant theme the model emits greens and teals everywhere as background or prop colour). These hues are removed from the pool for the matching theme to avoid the same scene-fill / outline confusion.
\end{enumerate}

\paragraph{Stage 2: outline rendering.}
The catalogue prompt is rewritten in \emph{outline mode}. The model is told, in unusually strict language, that it must draw only a four-pixel solid stroke in the assigned marker colour around each product's silhouette, with no fill, no gradient, no glow, no drop shadow and no surrounding decoration. Three further rules support detection at Stage 3.

\begin{itemize}
\item An eighty-pixel \emph{quiet zone} around each outline is reserved: no overlay text, price tag or decorative element may enter it. The quiet zone gives the deterministic mask in Stage 3 enough room to be dilated without erasing surrounding scene content.
\item Product orientation is constrained to that of the reference photograph: the model may rotate the product by a multiple of ninety degrees but is forbidden from inventing novel poses. Outlines drawn around invented poses are still detectable but bear no relation to the contour of the reference photograph that will be pasted in their place, which produces visible scale or aspect mismatches.
\item Drop shadows and reflections under the placeholder regions are explicitly forbidden, because they survive the outline-detection mask and produce ``ghost'' shadows beneath the pasted product after compositing.
\end{itemize}

The provider call is issued through the existing \texttt{IImageProvider} abstraction (Listing~\ref{lst:iimageprovider}), so that the choice of multimodal model is an orthogonal configuration decision rather than a property of the pipeline.

\paragraph{Stage 3: detect-and-composite.}
The third stage runs entirely on the application server, with no further model calls (other than the judge described in Section~\ref{subsec:hyb-judge}). It proceeds in six sub-steps:

\begin{enumerate}
\item[\textbf{(a)}] \textbf{Per-product colour matching.} Each pixel of the placeholder canvas is mapped to the closest marker colour by Chebyshev RGB distance under a tolerance of $80$, producing a per-pixel product index (or $-1$ if no marker is within tolerance).
\item[\textbf{(b)}] \textbf{Largest-component filter.} For each product, only the largest connected component of matched pixels is kept, discarding stray marker-coloured pixels in the background that would otherwise inflate the paste-back bounding box.
\item[\textbf{(c)}] \textbf{Outline thickening.} The retained outline is dilated by four pixels to absorb anti-aliased halo pixels at the stroke edge.
\item[\textbf{(d)}] \textbf{Inpaint mask construction.} A binary product-region mask is built by flood-filling unmatched cells from the canvas border and inverting the result, then dilated by fifteen pixels to also catch the model's occasional drop shadows.
\item[\textbf{(e)}] \textbf{LaMa inpainting.} The placeholder canvas and the mask are passed to a LaMa inpainting model \cite{suvorov2022lama} (served via IOPaint), chosen because it handles large masks, runs on CPU at acceptable latency, and needs no inference-time text prompt.
\item[\textbf{(f)}] \textbf{Product paste.} The retailer's source photograph is trimmed to its opaque bounding box, scaled by area to $0.8\times$ the median placeholder area, anchored at the placeholder's bottom-left corner so products ``stand'' on a baseline, and alpha-composited onto the inpainted canvas.
\end{enumerate}

Listing~\ref{lst:hybrid-loop} shows a slightly simplified version of the compositor's main loop; the production code performs the same steps with additional logging and per-product diagnostics that have been elided for legibility.

\begin{lstlisting}[style=csharp, caption={Per-product compositing loop (simplified).}, label={lst:hybrid-loop}]
foreach (var (product, marker) in assignments)
{
    var raw       = MatchByChebyshev(placeholder, marker, tolerance: 80);
    var component = LargestConnectedComponent(raw);
    var outline   = Dilate(component, radius: 4);
    var mask      = BuildProductMask(outline);
    mask          = Dilate(mask, radius: 15);

    var cleaned   = await lama.InpaintAsync(placeholder, mask, ct);
    var trimmed   = TrimTransparentBorder(product.Photo);
    var scaled    = ScaleByArea(trimmed, target: medianBboxArea * 0.8);
    var anchor    = BottomLeftOf(BoundingBox(component));
    placeholder   = AlphaComposite(cleaned, scaled, anchor);
}
return placeholder;
\end{lstlisting}

\subsection{The quality-gate judge and empirical observations}
\label{subsec:hyb-judge}

The Stage~3 output is a candidate final image, not the final image itself. A separate language-model call, the \emph{judge}, is given the candidate and asked, in one sentence, whether every product is acceptable: no leftover coloured outlines, no obvious paste seams, no major distortion of the pasted product. The judge is required to reply with a single word, either \texttt{YES} or \texttt{NO}. If the answer is \texttt{NO}, the entire pipeline (from Stage~2 onwards, with the marker assignment from Stage~1 preserved) is re-run. The loop is capped at three trials. If the judge service itself fails for any reason, the most recent candidate is accepted: the platform degrades open rather than blocking the retailer, on the principle that a marginally imperfect catalogue is more valuable than no catalogue at all.

The choice of an off-the-shelf multimodal model as judge, rather than a learned classifier, is itself part of the protocol-over-training argument. The judge's job is not to score image realism in general but to detect a small list of well-defined compositional faults (residual marker pixels, visible seams at the paste boundary, gross distortion of the pasted product). Off-the-shelf multimodal models are reliable on well-defined visual checklists of this shape, are cheap to call once per candidate, and adapt to new failure modes by a prompt edit rather than a retraining cycle.

\label{subsec:hyb-empirical}
The pipeline was exercised informally during development on a sample of catalogues drawn across several shop profiles, layouts and product counts. Two observations are worth recording.

The first is that the single-shot pipeline (Stages~1--3 followed by one judge call) was accepted by the judge in roughly seventy per cent of runs. The dominant failure mode at this rate was residual outline pixels at the boundary of an inpainted region, which the dilation in step~(d) of Stage~3 normally absorbs but occasionally misses when the model emits a partial outline that loops back on itself; the second-most-common failure mode was visible scale mismatch between a pasted product and its placeholder bounding box.

The second is that on a smaller sub-sample where the judge's verdict was compared against a human reviewer, the two agreed in nearly every case. The human reviewer was asked the same one-sentence question the judge was asked, and the disagreements were a small handful in which the judge was strict (rejecting an image the human was happy to accept) rather than the reverse. This is consistent with the observation, across the broader literature, that strong multimodal models are reliable critics of compositional faults that are easy to describe verbally, and is the platform's evidence that the judge is a good proxy for the human reviewer that it replaces in production.

If the per-trial acceptance probabilities are treated as approximately independent, capping the loop at three trials gives an effective acceptance rate of $1 - 0.30^{3} \approx 0.973$. Table~\ref{tab:hyb-acceptance} reports the corresponding rates for one through five trials, so that the diminishing marginal gain from raising the cap is visible at a glance. The pipeline could be extended to four trials ($\approx 0.992$) or five ($\approx 0.998$) with a strict mathematical improvement in acceptance probability, at the cost of additional latency and a proportional rise in the number of multimodal calls per delivered catalogue. The application keeps the cap at three because a $\approx 97\%$ acceptance rate is already adequate for the small-retailer workflow, and because the latency and cost added by a fourth or fifth trial fall on every user even though only the worst few per cent of runs would benefit; raising the cap remains a one-line change should the operational data later justify it.

\begin{table}[!ht]
\centering
\caption{Compound acceptance probability table.}
\label{tab:hyb-acceptance}
\small
\renewcommand{\arraystretch}{1.3}
\begin{tabularx}{\textwidth}{|>{\centering\arraybackslash}p{2.4cm}|>{\centering\arraybackslash}p{3.6cm}|Y|}
\hline
\textbf{Trials} & \textbf{Acceptance probability} & \textbf{Interpretation} \\
\hline
1 & $\approx 0.70$ & Single-shot baseline. \\
\hline
2 & $1 - 0.30^{2} = 0.91$ & At most one retry. \\
\hline
3 & $1 - 0.30^{3} \approx 0.97$ & Pipeline cap; the user-facing rate. \\
\hline
4 & $1 - 0.30^{4} \approx 0.992$ & Available if the cap were raised; not adopted. \\
\hline
5 & $1 - 0.30^{5} \approx 0.998$ & Diminishing return, omitted in production. \\
\hline
\end{tabularx}
\end{table}

These figures are observational and informal; a controlled benchmark with a fixed evaluation set, blind human raters and per-failure-mode breakdown is left to future work. They are presented here to motivate the three-trial cap, not as a definitive measurement.

\subsection{In-house composite onto a wallpaper}
\label{subsec:inhouse-composite}

The third catalogue flavour listed in Table~\ref{tab:catalogue-flavours} of Chapter~\ref{chap:architecture} is the in-house composite. It differs from the two flavours above in one decisive way: it invokes no generative model at the moment the catalogue is built. Where the fully generative mode (Section~\ref{subsec:gen-overview}) asks the provider to render the whole scene, and the preserve mode (Section~\ref{subsec:hyb-motivation}) asks it to render a scene with placeholder regions, the in-house composite asks the provider for nothing at all. The whole catalogue is assembled by the platform's own deterministic code.

The backdrop comes from one of two sources. The retailer can reuse a brand wallpaper produced earlier by the wallpaper service (Section~\ref{subsec:gen-overview}), the same brand-themed image shown in Figure~\ref{fig:result-wallpaper}, or upload a custom background of their own. Either way the backdrop is fixed before the catalogue is built, so the platform makes no inference call while assembling it.

Onto that backdrop the deterministic compositor lays the selected products. Each product photograph is trimmed to its opaque bounding box and alpha-composited into place, together with its price and the shop's brand identity, using the same compositing machinery that Stage~3 of the hybrid pipeline (Section~\ref{subsec:hyb-pipeline}) uses to paste real products into a generated scene. The difference is that here there is no generated scene to detect, no marker outline to erase and no judge loop to run, because nothing in the image was produced by a model and so nothing can be hallucinated or substituted.

The trade-off is straightforward. The in-house composite spends no credits at generation time, is fully deterministic and preserves every product pixel for pixel, which makes it the cheapest and most predictable of the three flavours. The price paid is that the visual richness of the result is bounded by the backdrop the retailer supplied or generated once, rather than tailored to the chosen products by a fresh model call. It therefore suits retailers who want complete control over the layout, or who want to produce catalogues repeatedly without spending credits on every one, while the two model-based flavours remain available whenever a freshly generated scene is worth the cost.

\section{Visual Product Search}
\label{sec:visual-search}

A retail asset is only as convincing as the product photograph it starts from. The platform therefore lets a retailer represent a product with a clean image from a curated catalogue of professional product photography, the evenly lit reference the hybrid catalogue pipeline of Section~\ref{subsec:hyb-pipeline} later reuses to keep products faithful. The theory behind this search by photo was covered in Section~\ref{sec:visual-similarity} and its storage in Section~\ref{subsec:relational-schema}. Figure~\ref{fig:visual-search} shows how the parts fit together.

The catalogue is curated in advance and organised by Category. Each image is embedded once by the CLIP ViT-B/32 vision encoder~\cite{radford2021clip}, which the platform hosts and runs on its own servers, and the resulting 512-dimensional vector is L2-normalised and stored in the \texttt{ReferenceImage.Embedding} column as a pgvector vector(512), indexed with a Hierarchical Navigable Small World index under cosine distance. Embedding happens ahead of time, so a search never re-runs the encoder over the catalogue.

When the retailer adds a product, they photograph the real item and the platform returns the ten closest catalogue products, ranked by visual similarity. They keep their own photo or pick one of these matches to represent the product.

\begin{figure}[H]
\centering
% --- self-contained palette (muted, print-friendly) ---
\definecolor{cInk}{HTML}{2E3440}    % text + arrows
\definecolor{cIndigo}{HTML}{3F4D8C} % LLM / orchestration stages
\definecolor{cSlate}{HTML}{4A5568}  % deterministic platform code
\definecolor{cTeal}{HTML}{2E5F8C}   % stores
\definecolor{cAmber}{HTML}{B07A1F}  % highlight accent
\definecolor{cGray}{HTML}{7C8390}   % external providers
\resizebox{\textwidth}{!}{%
\begin{tikzpicture}[
    font=\sffamily\scriptsize,
    card/.style={draw=#1, line width=0.9pt, rounded corners=3pt, fill=#1!7,
                 align=center, text=cInk,
                 drop shadow={shadow xshift=0.5mm, shadow yshift=-0.5mm,
                              fill=black, opacity=0.12}},
    group/.style={draw=#1, line width=0.9pt, rounded corners=3pt, fill=#1!6,
                  inner sep=2.8mm,
                  drop shadow={shadow xshift=0.5mm, shadow yshift=-0.5mm,
                               fill=black, opacity=0.12}},
    tab/.style={fill=#1, rounded corners=2.5pt, inner xsep=6pt, inner ysep=0pt,
                minimum height=5.2mm, anchor=south west,
                font=\sffamily\scriptsize\bfseries, text=white},
    chip/.style={draw=#1, line width=0.6pt, rounded corners=2.5pt, fill=white,
                 align=center, text=cInk},
    store/.style={draw=cTeal, line width=0.6pt, cylinder, shape border rotate=90,
                  aspect=0.22, minimum width=24mm, minimum height=15mm, align=center,
                  cylinder uses custom fill, cylinder body fill=white,
                  cylinder end fill=cTeal!15, text=cInk},
    arr/.style={-{Latex[length=2mm,width=1.5mm]}, line width=0.7pt, draw=cInk!85},
    fb/.style={-{Latex[length=2mm,width=1.5mm]}, line width=0.7pt, draw=cGray,
               dash pattern=on 2.4pt off 1.7pt},
    arrlbl/.style={font=\sffamily\tiny\itshape, fill=white, inner sep=1pt, text=cInk!80}
]

% --- inputs ---
\node[store] (cat) at (-4.2, 1.7) {Curated\\catalogue};
\node[card=cTeal, minimum width=20mm, minimum height=11mm] (photo) at (-4.2, -1.7)
    {phone\\photo};

% --- shared encoder (hero) and the vector it emits ---
\node[card=cIndigo, minimum width=32mm, minimum height=16mm] (enc) at (0, 0)
    {CLIP \texttt{ViT-B/32}\\vision encoder\\(ONNX)};
\node[chip=cSlate, minimum width=26mm, minimum height=12mm] (vec) at (3.4, 0)
    {512-d vector\\(L2-normalised)};

% --- offline destination ---
\node[store] (pg) at (7.4, 1.7) {pgvector index\\\texttt{vector(512)}\\HNSW, cosine};

% --- online query path ---
\node[card=cSlate, minimum width=30mm, minimum height=13mm] (topk) at (7.2, -1.7)
    {cosine top-10\\nearest neighbours};
\node[card=cAmber, minimum width=22mm, minimum height=13mm] (cand) at (10.8, -1.7)
    {ranked\\matches};
\node[store] (lib) at (13.8, -1.7) {product\\library};

% --- central frame: the one shared encoder ---
\begin{pgfonlayer}{background}
\node[group=cIndigo, fit=(enc)(vec)] (encbox) {};
\end{pgfonlayer}
\node[tab=cIndigo] at ([xshift=4mm]encbox.north west) {shared CLIP encoder};

% --- arrows ---
\draw[fb]  (cat.east)   to[out=0, in=150] node[arrlbl, pos=0.55]{once, on import} (enc.150);
\draw[arr] (photo.east) to[out=0, in=210] node[arrlbl, pos=0.55]{on upload} (enc.210);
\draw[arr] (enc.east) -- (vec.west);
\draw[fb]  (vec.north) to[out=90, in=180]  node[arrlbl, pos=0.62]{index} (pg.west);
\draw[arr] (vec.south) to[out=-90, in=180] node[arrlbl, pos=0.62]{query} (topk.west);
\draw[fb]  (pg.south) to[out=-90, in=90] node[arrlbl, pos=0.5]{indexed vectors} (topk.north);
\draw[arr] (topk.east) -- node[arrlbl]{top 10} (cand.west);
\draw[arr] (cand.east) -- node[arrlbl]{keep or swap} (lib.west);

\end{tikzpicture}%
}
\caption{Search by photo. The curated catalogue is embedded once by the CLIP vision encoder and indexed in pgvector. An uploaded phone photograph is embedded by the same encoder and matched by cosine nearest-neighbour search, and the retailer keeps the photograph or swaps in the closest studio match.}
\label{fig:visual-search}
\end{figure}

\section{Print Pipeline and DPI Enhancer}
\label{sec:print-pipeline}

The print pipeline is the third place where a deterministic post-processor closes the gap between what the multimodal model emits and what the retailer needs. Generated assets ship at the model's native resolution, typically $1024 \times 1024$ pixels. A print of the same asset at $300\,$dpi requires materially more pixels: $1240 \times 1748$ for A6, rising to $3508 \times 4961$ for A3. The print pipeline therefore upscales below-target assets by a learned super-resolution model before laying them out in a paper-sized PDF.

The model used is Real-ESRGAN \cite{wang2021realesrgan}, exported to ONNX and run through the application server's ONNX Runtime. Inputs at most $512$ pixels on a side are upscaled in a single forward pass; larger inputs are tiled into $256 \times 256$ source patches with sixteen pixels of overlap padding on each side, each patch is upscaled independently, and the upscaled patches are stitched back together. The overlap padding is what makes the tiling invisible: each interior pixel of the output appears in two adjacent tiles' interiors, and the stitching averages them, which suppresses the seam artefacts that naively tiled super-resolution produces. Tile-based inference also caps the model's peak memory consumption at the size of a single patch's intermediate activations, which keeps the upscaler runnable inside the same Container Apps revision as the rest of the API. Concurrent inference is bounded by a small semaphore, so that an unusually large print job cannot crowd out interactive traffic. Figure~\ref{fig:print-flow} sketches the print flow end to end.

\begin{figure}[H]
\centering
% --- self-contained palette (muted, print-friendly) ---
\definecolor{cInk}{HTML}{2E3440}    % text + arrows
\definecolor{cIndigo}{HTML}{3F4D8C} % LLM / orchestration stages
\definecolor{cSlate}{HTML}{4A5568}  % deterministic platform code
\definecolor{cTeal}{HTML}{2E5F8C}   % stores
\definecolor{cAmber}{HTML}{B07A1F}  % highlight accent
\definecolor{cGray}{HTML}{7C8390}   % external providers
\resizebox{\textwidth}{!}{%
\begin{tikzpicture}[
    font=\sffamily\scriptsize,
    inp/.style={draw=cTeal, line width=0.9pt, rounded corners=3pt, fill=cTeal!7,
                minimum width=30mm, minimum height=12mm, align=center, text=cInk,
                font=\sffamily\scriptsize\itshape,
                drop shadow={shadow xshift=0.5mm, shadow yshift=-0.5mm,
                             fill=black, opacity=0.12}},
    proc/.style={draw=cSlate, line width=0.9pt, rounded corners=3pt, fill=cSlate!7,
                 minimum width=32mm, minimum height=12mm, align=center, text=cInk,
                 drop shadow={shadow xshift=0.5mm, shadow yshift=-0.5mm,
                              fill=black, opacity=0.12}},
    aicall/.style={draw=cGray, line width=0.9pt, rounded corners=3pt, fill=cGray!6,
                   dash pattern=on 2.4pt off 1.7pt, minimum width=32mm,
                   minimum height=15mm, align=center, text=cInk},
    final/.style={draw=cTeal, line width=0.6pt, cylinder, shape border rotate=90,
                  aspect=0.22, minimum width=26mm, minimum height=14mm, align=center,
                  cylinder uses custom fill, cylinder body fill=white,
                  cylinder end fill=cTeal!15, text=cInk},
    group/.style={draw=#1, line width=0.9pt, rounded corners=3pt, fill=#1!6,
                  inner sep=2.8mm,
                  drop shadow={shadow xshift=0.5mm, shadow yshift=-0.5mm,
                               fill=black, opacity=0.12}},
    tab/.style={fill=#1, rounded corners=2.5pt, inner xsep=6pt, inner ysep=0pt,
                minimum height=5.2mm, anchor=south west,
                font=\sffamily\scriptsize\bfseries, text=white},
    arr/.style={-{Latex[length=2mm,width=1.5mm]}, line width=0.7pt, draw=cInk!85},
    barr/.style={-{Latex[length=2mm,width=1.5mm]}, line width=0.7pt, draw=cInk!70,
                 dash pattern=on 2.4pt off 1.7pt},
    arrlbl/.style={font=\sffamily\tiny\itshape, fill=white, inner sep=1pt,
                   text=cInk!80}
]

% --- gallery endpoints (outside the pipeline frame) ---
\node[inp]    (src) at (0, 0)     {gallery image\\(native resolution)};
\node[proc]   (chk) at (4.6, 0)   {compute target pixels\\(paper size @ $300\,$dpi)};
\node[aicall] (up)  at (9.4, 1.9) {Real-ESRGAN\\tile upscaler\\[1pt]\tiny\textsc{ai call (ONNX)}};
\node[proc]   (lay) at (14.2, 0)  {PDF layout\\($\pm$ QR code)};
\node[final]  (out) at (18.6, 0)  {print-ready\\PDF in gallery};

% --- frame the processing stages (gallery endpoints stay outside) ---
\begin{pgfonlayer}{background}
\node[group=cSlate, fit=(chk)(up)(lay)] (pipe) {};
\end{pgfonlayer}
\node[tab=cSlate] at ([xshift=4mm]pipe.north west) {Print pipeline};

% --- flow with a conditional upscale gate ---
\draw[arr] (src) -- (chk);
\fill[cInk!75] (chk.east) circle (1.4pt);
\draw[arr]  (chk.east) to[out=55, in=180]
      node[arrlbl, pos=0.55] {if scale~$>1$} (up.west);
\draw[arr]  (up.east) to[out=0, in=95] (lay.north);
\draw[barr] (chk.east) -- node[arrlbl] {else} (lay.west);
\draw[arr]  (lay) -- (out);

\end{tikzpicture}%
}
\caption{Print pipeline: assets are upscaled only when the target paper size at $300\,$dpi demands more pixels than the asset carries.}
\label{fig:print-flow}
\end{figure}

